% ==============================================================================
% PG - Lorenzo Rizzi Fiorot
% Capítulo 4 - O EduGraph
% ==============================================================================
% TODO(autor): o modelo do curso sugere que o título deste capítulo nomeie a
% contribuição do trabalho; "O EduGraph" é provisório.
\chapter{O EduGraph}
\label{sec-edugraph}

Texto.


%%% Início de seção. %%%
\section{Requisitos}
\label{sec-edugraph-requisitos}

% TODO(autor): os requisitos funcionais e não funcionais foram registrados em algum
% documento? Sem esse registro, a seção não descreve requisitos que não existam.
Texto.


%%% Início de seção. %%%
\section{Arquitetura}
\label{sec-edugraph-arquitetura}

O EduGraph organiza-se em três camadas, ilustradas na Figura~\ref{fig-arquitetura}: a interface web, executada no navegador; a API, que concentra o acesso aos dados; e o banco de grafos Neo4j, que armazena o repositório de conhecimento. Para obter os dados, o navegador comunica-se apenas com o endereço do próprio sistema: em desenvolvimento, o servidor do Vite entrega a interface e repassa à API os pedidos cujo caminho começa com \texttt{/api/}, de modo que a interface não conhece o endereço do banco nem as credenciais de acesso. Na versão atual, todos esses serviços são executados localmente; a publicação do sistema depende da definição da hospedagem.

% Diagrama feito só com caixas e setas do próprio LaTeX, sem pacotes de desenho, para não
% aumentar o tempo de compilação.
\begin{figure}[htb]
    \centering
    \caption{Arquitetura do EduGraph}
    \label{fig-arquitetura}
    \setlength{\fboxsep}{7pt}
    \renewcommand{\arraystretch}{1.2}
    \begin{tabular}{c}
        \fbox{\parbox{12cm}{\centering\textbf{Navegador}\\[2pt]
            Interface web em React: páginas, busca guiada\\
            e visualização do grafo em Canvas 2D}} \\
        $\Big\updownarrow$\ {\small HTTP: a interface e as rotas \texttt{/api/}, com respostas em JSON} \\
        \fbox{\parbox{12cm}{\centering\textbf{Servidor do Vite} (em desenvolvimento)\\[2pt]
            Entrega a interface e repassa à API os pedidos que começam com \texttt{/api/}}} \\
        $\Big\updownarrow$\ {\small HTTP: rotas da API, com respostas em JSON} \\
        \fbox{\parbox{12cm}{\centering\textbf{API em Python (FastAPI)}\\[2pt]
            Consultas da aplicação por assunto e verificação do serviço;\\
            guarda a credencial do banco e impõe o modo somente leitura;\\
            módulos-base do fórum, das ontologias e do assistente}} \\
        $\Big\updownarrow$\ {\small Query API (HTTP): Cypher em modo de leitura, resposta em JSON tipado} \\
        \fbox{\parbox{12cm}{\centering\textbf{Banco de grafos Neo4j} (edição Community)\\[2pt]
            Repositório de conhecimento: o grafo da BNCC (Seção~\ref{sec-edugraph-modelo})}} \\
    \end{tabular}
\end{figure}

O contrato entre a interface e a API baseia-se em requisições HTTP e respostas em JSON. Cada consulta da aplicação --- os números da página inicial, os anos, as áreas do conhecimento e os componentes de um ano, o recorte de um ano em um componente, as perguntas da busca guiada e as conexões de um nó --- é exposta por uma rota da API que executa uma única instrução Cypher, sempre com parâmetros, de modo que os valores informados nunca são inseridos no texto da consulta. A resposta reúne a instrução executada, os parâmetros e o resultado devolvido pelo Neo4j, sem alteração, no formato JSON tipado; com isso, a interface reconstrói os nós, as relações e os caminhos de qualquer resultado e exibe ao usuário a consulta realizada (Seção~\ref{sec-edugraph-grafos}). Os erros do banco, da API e do servidor de desenvolvimento seguem um formato único, com um código e uma mensagem, o que permite à interface tratá-los de forma uniforme.

A credencial de acesso ao banco fica apenas na configuração da API, em um arquivo local excluído do controle de versão, e nunca é enviada ao navegador. Toda consulta é encaminhada ao Neo4j em modo de leitura, independentemente do que a interface solicite, e o próprio servidor recusa instruções de escrita nesse modo, comportamento verificado durante o desenvolvimento. A medida é necessária porque o controle de acesso por papéis é recurso da edição Enterprise do Neo4j \cite{neo4j2026autorizacao}: na edição Community, adotada no projeto, não há como restringir um usuário à leitura. Por fim, a API valida os parâmetros recebidos, como os códigos de ano, de área e de componente e os limites de resultados, e limita a 12 segundos a espera por uma resposta do banco.

A API reúne, ainda, uma rota que informa se o serviço e a conexão com o banco estão disponíveis e os módulos-base de três serviços previstos --- o fórum, o repositório de ontologias e o assistente ---, cujas rotas já constam da documentação interativa gerada pelo FastAPI e respondem como ainda não implementadas (Seção~\ref{sec-edugraph-api}).


%%% Início de seção. %%%
\section{Modelo do Grafo}
\label{sec-edugraph-modelo}

% TODO(autor): dizer aqui quem definiu o modelo do grafo e quem mantém a base; a redação
% depende da divisão de responsabilidades descrita na Introdução.
O repositório de conhecimento é um grafo de propriedades armazenado no Neo4j, cujo modelo reproduz o encadeamento da BNCC descrito na Seção~\ref{sec-referencial-bncc}: ano, área do conhecimento, componente curricular, unidade temática, objeto de conhecimento e habilidade. Cada nível corresponde a nós com rótulo próprio e liga-se ao nível seguinte por uma relação nomeada, conforme a Tabela~\ref{tab:modelo-grafo}; as relações não têm propriedades. Nos três últimos níveis, o rótulo leva o sufixo do componente curricular --- \texttt{Mat}, para Matemática, e \texttt{Cien}, para Ciências ---, como em \texttt{HabilidadeMat}. Além do encadeamento, cada ano liga-se diretamente às suas habilidades pela relação \texttt{DESENVOLVE\_HABILIDADE}.

% TODO(autor): na base atual, as descrições das habilidades de Ciências coincidem com o
% texto da BNCC, e as de Matemática estão resumidas. Se continuarem assim, dizer isso aqui.
\begin{table}[h]
    \centering
    \caption{Níveis do modelo do grafo}
    \label{tab:modelo-grafo}
    \renewcommand{\arraystretch}{1.3}
    \begin{tabular}{|>{\raggedright\arraybackslash}p{4.2cm}|>{\raggedright\arraybackslash}p{5.6cm}|>{\raggedright\arraybackslash}p{4.6cm}|}
        \hline
        \textbf{Nível} & \textbf{Rótulo dos nós} & \textbf{Relação com o nível seguinte} \\
        \hline
        Ano & \texttt{Ano} & \texttt{TEM\_AREA} \\
        \hline
        Área do conhecimento & \texttt{AreaConhecimento} & \texttt{TEM\_COMPONENTE} \\
        \hline
        Componente curricular & \texttt{ComponenteCurricular} & \texttt{TEM\_UNIDADE} \\
        \hline
        Unidade temática & \texttt{UnidadeTematica} + sufixo & \texttt{TEM\_OBJETO} \\
        \hline
        Objeto de conhecimento & \texttt{ObjetoConhecimento} + sufixo & \texttt{DESENVOLVE} \\
        \hline
        Habilidade & \texttt{Habilidade} + sufixo & --- \\
        \hline
    \end{tabular}
\end{table}

Os nós de área, de componente, de unidade temática e de objeto de conhecimento não se repetem por ano: um nó é compartilhado por todos os anos em que aparece, como a unidade temática Terra e Universo, presente nos quadros do 1º e do 2º ano, à qual se ligam os objetos de conhecimento dos dois anos. O que associa cada habilidade a um ano é a relação direta \texttt{DESENVOLVE\_HABILIDADE}; as consultas da aplicação partem dela para restringir os resultados ao ano escolhido, o que impede que conteúdos de anos distintos se misturem nos subgrafos exibidos. Cada ano, por sua vez, liga-se apenas às áreas do conhecimento em que tem conteúdo.

Embora os rótulos variem conforme o componente, as consultas com que a aplicação obtém os dados percorrem o grafo pelas relações, comuns a todos os componentes, e só mencionam os rótulos \texttt{Ano}, \texttt{AreaConhecimento} e \texttt{ComponenteCurricular}. Assim, um componente carregado com a mesma estrutura passa a ser consultado sem alteração do código.

Em 26 de novembro de 2026, a base continha 75 nós e 108 relações, distribuídos conforme a Tabela~\ref{tab:cobertura-grafo}. Dos quatro nós de área do conhecimento, dois --- Ciências da Natureza e Matemática --- ligam-se a componentes curriculares; Linguagens e Ciências Humanas ainda não têm componentes carregados.

\begin{table}[h]
    \centering
    \caption{Conteúdo da base de dados em 26 de novembro de 2026}
    \label{tab:cobertura-grafo}
    \renewcommand{\arraystretch}{1.3}
    \begin{tabular}{|l|l|c|c|c|}
        \hline
        \textbf{Ano} & \textbf{Componente} & \textbf{Unidades temáticas} & \textbf{Objetos} & \textbf{Habilidades} \\
        \hline
        1º ano & Matemática & 5 & 17 & 22 \\
        \hline
        1º ano & Ciências & 3 & 3 & 6 \\
        \hline
        2º ano & Ciências & 3 & 3 & 8 \\
        \hline
        \multicolumn{2}{|l|}{Nós distintos} & 8 & 23 & 36 \\
        \hline
    \end{tabular}
\end{table}

\noindent As três unidades temáticas de Ciências são as mesmas nos dois anos, o que explica o total de oito unidades distintas. Somam-se aos nós da tabela os 2 nós de ano, os 4 de área do conhecimento e os 2 de componente curricular. Das 108 relações, 36 são do tipo \texttt{DESENVOLVE\_HABILIDADE}, uma para cada habilidade.

% TODO(autor): a base não tem nós de conceitos do Pensamento Computacional nem habilidades
% da BNCC Computação; dizer aqui como o trabalho trata isso, conforme o enquadramento
% definido na Introdução.


%%% Início de seção. %%%
\section{Prototipação}
\label{sec-edugraph-prototipacao}

Texto.


%%% Início de seção. %%%
\section{Interface e Identidade Visual}
\label{sec-edugraph-interface}

Texto.


%%% Início de seção. %%%
\section{Página Inicial}
\label{sec-edugraph-inicial}

Texto.


%%% Início de seção. %%%
\section{Página de Grafos}
\label{sec-edugraph-grafos}

Texto.


%%% Início de seção. %%%
\section{Acessibilidade}
\label{sec-edugraph-acessibilidade}

Texto.


%%% Início de seção. %%%
\section{API e Base de Dados}
\label{sec-edugraph-api}

Texto.


%%% Início de seção. %%%
\section{Telas Complementares}
\label{sec-edugraph-complementares}

% TODO(autor): as telas de conta, o Fórum e as Ontologias entram aqui como telas
% implementadas só no front-end, com conteúdo ilustrativo, ou só como trabalho futuro?
Texto.
